% Sec. VI: Bypass resistance for routes of any length (VI-A to VI-D: lightning.spthy),
% VI-E: the PTLC countermeasure (lightning_aux.spthy -D=PTLC / -D=HTLC).
\section{Bypass Resistance for Routes of Any Length}\label{sec:bypass}

The wormhole of Sec.~\ref{sec:wormhole} lets a payment settle while an honest forwarder between
two compromised parties is bypassed: its outgoing HTLC is refunded and it earns no fee. We now
show that in our model a bypass requires key compromise: for the forwarder at each of the
positions 1 to 6 of a route, and, for payments whose hash is offered on a single channel, for a
forwarder at any position of a route of any length. Secs.~\ref{sec:bypass-fixed}
to~\ref{sec:bypass-caveats} concern the main theory \code{lightning.spthy}; there is no separate
N-hop theory.

HTLC offers between channel peers are not network messages in our model. When tracing a
forwarded message, \tamarin{} treats every party as a potential sender, and with the offer as a
public, signed message, the analysis of routes of unbounded length did not terminate. We
therefore swapped the public message for a private, single-use fact, \code{Route}, which each
application of the one generic rule \code{Forward_HTLC} consumes and re-creates for the next hop
(Sec.~\ref{sec:core-routing}); what this takes away from the adversary is stated in
Sec.~\ref{sec:limitations}.

\subsection{Bypass Requires Compromise}\label{sec:bypass-fixed}

The only rule that produces a \code{Fulfilled} fact for an HTLC that was not redeemed is
\code{AttackerSkip}. It requires the two parties at the ends of the skip, \code{ExitAtt} and
\code{EntryAtt}, to be compromised, and it records \code{Skipped(ptrSkip,x)};
\code{Fulfilled_Source} states that every \code{Fulfilled} fact used in a settlement stems from
an earlier redemption or skip. Every skip follows a compromise (Lemma~\ref{lem:skip}).

A forwarder is \emph{bypassed} when the payment settles at the sender while the forwarder's own
hop is refunded; \code{Intermediary_Timeout_Refund} records such a refund as
\code{RefundHop(ptrIn,ptrOut,y)}. For the forwarders at the first six positions of the sender's
route we prove, in the full model with T1, T2 and T3 in place:

\begin{theorem}[Bypass Resistance at Positions 1 to 6]
Let $1 \le k \le 6$. If a sender settles a payment with preimage $x$ on its own channel, the
HTLC on $h(x)$ is forwarded hop by hop from that channel by forwarders $F_1, \dots, F_k$, and
the hop of $F_k$ is refunded, then some party's long-term key was compromised.
\end{theorem}
For $k = 1$ and $k = 2$, the statements are:
\begin{spthy}
lemma First_Hop_Skip_Requires_Compromise [heuristic={lemmas_first}]:
    "All S F1 R pSF Next pOut x #s #f #r.
        SenderSettled(S, F1, R, pSF, x) @ #s
      & HTLCForwarded(F1, S, Next, pSF, pOut, h(x)) @ #f
      & RefundHop(pSF, pOut, h(x)) @ #r
     ==> Ex P #c. LtkCompromised(P) @ #c"
\end{spthy}
\begin{spthy}
lemma Second_Hop_Skip_Requires_Compromise [heuristic={lemmas_first}]:
    "All S F1 R pSF x F2 p1 Next pOut #s #f1 #f2 #r.
        SenderSettled(S, F1, R, pSF, x) @ #s
      & HTLCForwarded(F1, S, F2, pSF, p1, h(x)) @ #f1
      & HTLCForwarded(F2, F1, Next, p1, pOut, h(x)) @ #f2
      & RefundHop(p1, pOut, h(x)) @ #r
     ==> Ex P #c. LtkCompromised(P) @ #c"
\end{spthy}
For the forwarder at positions 3 to 6, the lemmas \code{Hop3_Skip_Requires_Compromise} to
\code{Hop6_Skip_Requires_Compromise} state the same, with the route from the sender's channel
spelled out as a chain of \code{HTLCForwarded} facts (Appendix~\ref{app:catalogue}). These
six theorems and their seven helper lemmas (group 9 of Table~\ref{tab:groups}), the 13
hop-bypass lemmas of \texttt{make verify}, are verified in the full model, with all
restrictions, and need no assumption on how often a payment hash is offered.

In the wormhole trace \code{Wormhole_Fee_Theft_Reachable} (Fig.~\ref{fig:wormhole}), Indiana is
the second forwarder (\code{F1} is Att1, \code{F2} is Indiana): Alice settles while Indiana's
outgoing HTLC is refunded, and Alice, Bob and Indiana are uncompromised. The trace satisfies the
premise of \code{Second_Hop_Skip_Requires_Compromise}, and the compromise in the conclusion is
that of Att1 and Att2, which \code{AttackerSkip} requires: the conclusion cannot be dropped.

\subsection{Routes of Any Length}\label{sec:bypass-any}

Each fixed-position theorem names the route up to the bypassed forwarder and therefore covers
one position only. The general theorem quantifies over any forward of the payment hash instead.

\begin{theorem}[Bypass Resistance at Any Position]\label{thm:anyhop}
If a sender settles a payment with preimage $x$ on its channel $p_{SF}$, every HTLC on $h(x)$
is offered on $p_{SF}$, and some forwarder of an HTLC on $h(x)$, at any position, has its hop
refunded, then some party's long-term key was compromised.
\end{theorem}
\begin{spthy}
lemma Any_Hop_Skip_Requires_Compromise [heuristic={lemmas_first}]:
  all-traces
  "∀ S F1 R pSF x Me Prev Next pIn pOut #s #f #r.
    ((((SenderSettled( S, F1, R, pSF, x ) @ #s) ∧
       (HTLCForwarded( Me, Prev, Next, pIn, pOut, h(x) ) @ #f)) ∧
      (RefundHop( pIn, pOut, h(x) ) @ #r)) ∧
     (∀ A B p #o. (HTLCOffered( A, B, p, h(x) ) @ #o) ⇒ (p = pSF))) ⇒
    (∃ P #c. LtkCompromised( P ) @ #c)"
\end{spthy}
The last conjunct of the premise is the \emph{single-offer premise}: every HTLC offered for
$h(x)$ is offered on the sender's channel \code{pSF}. No position appears in the statement:
\code{Me} is any forwarder of $h(x)$, at any distance from the sender.

\subsection{Proof}\label{sec:bypass-proof}

The proof is organised in seven helper lemmas, proved in the same run.
\code{Forward_Parent}: the incoming HTLC of every forward was offered by a sender or forwarded
by the previous hop. \code{Root_First}: every forward of a payment hash is, or is preceded by, a
forward whose incoming HTLC was offered by a sender. \code{Next_Is_Child}: if all offers of the
hash are on one channel and \code{Prev} forwarded an HTLC to \code{Me}, every later forward of
the hash is \code{Me}'s forward of that HTLC or comes after it. \code{Hop_Redeem_Propagates}:
without compromise, if a forwarder's incoming HTLC is redeemed, so is its outgoing one.
\code{Sender_Channel_Redeemed}: without compromise, if the sender settles, its own channel was
redeemed. \code{Route_Redeemed_Gen}: without compromise and if all offers of the hash are on one
channel, redemption of the incoming HTLC of a forward carries over to the incoming HTLC of every
forward of the hash that does not come earlier. \code{Route_Redeemed_No_Compromise} combines it
with \code{Sender_Channel_Redeemed} and \code{Root_First}: without compromise and under the
single-offer premise, if the sender settles, the incoming HTLC of every forward of the hash is
redeemed. \code{Root_First}, \code{Next_Is_Child} and \code{Route_Redeemed_Gen} are proved by
induction over traces.

The theorem then takes 7 proof steps. Assume no compromise. By
\code{Route_Redeemed_No_Compromise}, the bypassed forwarder's incoming HTLC \code{pIn} is
redeemed. By \code{Redeemed_Source}, this redemption is either the settlement of a forward from
\code{pIn}, which \code{Forward_Unique_Per_Incoming}, \code{SettledHop_From_Forward} and
\code{No_Settle_And_Refund_Same_Hop} exclude given \code{RefundHop(pIn,pOut,h(x))}, or a
fulfilment by the receiver, which \code{No_Fulfill_On_Forwarded_Incoming} excludes. The proof
scripts of \code{Root_First}, \code{Next_Is_Child}, \code{Route_Redeemed_Gen},
\code{Route_Redeemed_No_Compromise} and of the theorem are stored in \code{lightning.spthy}; the
other helpers are proved automatically.

\paragraph{General Mode and Trace Inclusion}
The theorem is proved in a separate mode of the same theory, \texttt{make general}
(\code{-D=GENERAL} \code{-D=NO_IMC} \code{-D=NO_OTF} \code{-D=NO_RBT} \code{-D=NECESSITY}),
which enables the general lemmas and removes the three liveness restrictions T1
(\code{RedeemBeforeTimeout}), T2 (\code{OutgoingTimesOutFirst}) and T3
(\code{IntermediaryMustClaim}). Removing a restriction only adds traces, so every trace of the
full model is a trace of this mode, and an all-traces property proved in this mode also holds in
the full model. The run proves the theorem together with all 23 \code{[reuse]} lemmas available
in this mode (24/24 verified), so no lemma proved under the removed restrictions is reused;
\code{ForwardTimeGap}, which is derived from T2, is not a \code{[reuse]} lemma in this mode.
Bypass resistance for single-offer payments therefore does not depend on T1, T2 or T3.

\subsection{Caveats}\label{sec:bypass-caveats}

\begin{itemize}
\item \emph{Single-offer premise.} Theorem~\ref{thm:anyhop} is proved only for payments whose
  hash is offered on a single channel. If the same hash is offered on several channels, several
  routes carry it, and the theorem says nothing about a forwarder on such a payment. The
  fixed-position theorems do not need the premise, because they spell out the route from the
  settling sender's channel; for such payments they cover positions 1 to 6 of that route.
\item \emph{Weak conclusion.} The conclusion states only that some party's long-term key was
  compromised at some point of the trace. It does not say which party (for instance a neighbour
  of the bypassed forwarder, as in the wormhole) and does not order the compromise before the
  settlement or the refund. The fixed-position theorems have the same conclusion; only
  \code{Skipped_Requires_Compromise} orders the compromise before the skip.
\item \emph{Liveness restrictions.} The six fixed-position theorems are proved in the full model,
  with T1, T2 and T3 in place, and we have not checked them without these restrictions. For a
  payment whose hash is offered on several channels, bypass resistance at positions 1 to 6
  therefore rests on the liveness restrictions; only Theorem~\ref{thm:anyhop} (single-offer
  payments) is shown to hold without them.
\end{itemize}

\subsection{Countermeasure: PTLC}\label{sec:ptlc}

The wormhole works because every channel of an HTLC route is locked with the same hash $h(x)$:
once $x$ is known at one point of the route, a compromised party can use it on a channel further
upstream. Malavolta et al.~\cite{malavolta2019} replace the shared lock by anonymous multi-hop
locks, the basis of point time-locked contracts (PTLCs): every channel has its own lock, and a
forwarder can compute the secret of its incoming lock only from the secret of its outgoing lock
and its own tweak, chosen by the sender. We compare both designs in two small models of
\code{lightning_aux.spthy} that differ only in how the channel secrets are derived and,
accordingly, in how an honest hop obtains its incoming secret (\code{Hop_Settle}):
\code{-D=PTLC} (per-hop secrets) and \code{-D=HTLC} (one shared secret, the baseline).

\paragraph{Model}
Both models fix the route S$\,\to\,$H1$\,\to\,$H2$\,\to\,$H3$\,\to\,$R over channels
\code{c0}, \code{c1}, \code{c2}, \code{c3}. \code{Sender_Setup} draws a fresh secret $x$ and
per-hop tweaks $t_1, t_2, t_3$ and derives the channel secrets by a hash chain
$s_{\mathit{in}} = h(\langle s_{\mathit{out}}, t\rangle)$: $s_3 = x$,
$s_2 = h(\langle s_3, t_3\rangle)$, $s_1 = h(\langle s_2, t_2\rangle)$,
$s_0 = h(\langle s_1, t_1\rangle)$; channel $c_i$ is locked with $h(s_i)$. An honest hop that
learns the secret of its outgoing channel derives the secret of its incoming channel with its
tweak (\code{Hop_Settle}). The hash chain is a symbolic stand-in for the additive tweak of PTLCs
built from adaptor signatures~\cite{malavolta2019}, $s_{\mathit{in}} = s_{\mathit{out}} + t$: in
both, deriving a hop's incoming secret requires that hop's tweak. In the HTLC baseline,
$s_0 = s_1 = s_2 = s_3 = x$, i.e., one hash lock $h(x)$ on every channel, and \code{Hop_Settle}
passes the outgoing secret on unchanged. The two rules of the PTLC model that implement the
per-hop locks are:
\begin{spthy}
rule Sender_Setup:
    let s3 = ~x
        s2 = h(<s3, ~t3>)
        s1 = h(<s2, ~t2>)
        s0 = h(<s1, ~t1>)
    in
    [ Fr(~x), Fr(~c0), Fr(~c1), Fr(~c2), Fr(~c3), Fr(~t1), Fr(~t2), Fr(~t3) ]
    --[ Pay($S, $H1, $H2, $H3, $R, ~c0, ~c1, ~c2, ~c3),
        Neq($S, $H1), Neq($S, $H2), Neq($S, $H3), Neq($S, $R),
        Neq($H1, $H2), Neq($H2, $H3), Neq($H1, $H3),
        Neq($H1, $R), Neq($H2, $R), Neq($H3, $R) ]->
    [ !Lock(~c0, $S, $H1, h(s0)), Open(~c0), SenderWait($S, ~c0),
      !Hop($H1, ~c0, ~c1, $H2, h(s1), ~t1),
      !Hop($H2, ~c1, ~c2, $H3, h(s2), ~t2),
      !Hop($H3, ~c2, ~c3, $R,  h(s3), ~t3),
      !RecvSecret($R, ~c3, ~x),
      Out(h(s0)), Out(h(s1)), Out(h(s2)), Out(h(s3)) ]
\end{spthy}
\begin{spthy}
rule Hop_Settle:
    [ Pend($H, cIn, cOut, t), ClaimInfo(cOut, sec), !Lock(cIn, P, $H, Lin), Open(cIn) ]
    --[ Redeemed(cIn), HopSettled($H, cIn, cOut), Eq(h(h(<sec, t>)), Lin) ]->
    [ ClaimInfo(cIn, h(<sec, t>)), Out(h(<sec, t>)) ]
\end{spthy}
The adversary is the same in both models: Dolev-Yao plus compromise of any party at any time
(rule \code{Compromise}). A compromised hop leaks its tweak (\code{Leak_Tweak}); every secret
used to claim a channel is published, as it would be on chain; a compromised payee may claim its
channel with any value the adversary knows and that opens the lock (\code{Adversary_Claim}).
There is no skip rule: a wormhole has to be built from adversary knowledge alone. Timing is
abstracted: an open HTLC may be refunded at any time (\code{Refund}).

\paragraph{Results}
With per-hop locks, the sender's channel \code{c0} is never redeemed while H2's outgoing channel
\code{c2} is refunded, as long as S and H2 are uncompromised, whichever of H1, H3 and R are
compromised (all-traces, verified, 25 steps):
\begin{spthy}
lemma No_Wormhole_Around_Honest_Hop:
  all-traces
  "not (
Ex S H1 H2 H3 R c0 c1 c2 c3 #p #a #b.
      Pay(S, H1, H2, H3, R, c0, c1, c2, c3) @ #p
    & Redeemed(c0) @ #a
    & Refunded(c2) @ #b
    & not (Ex #k. Compromised(H2) @ #k)
    & not (Ex #k. Compromised(S) @ #k))"
\end{spthy}
With the shared hash lock, the same situation is reachable (exists-trace, verified, 12 steps):
\begin{spthy}
lemma Wormhole_Around_Honest_Hop:
  exists-trace
  "Ex S H1 H2 H3 R c0 c1 c2 c3 #p #a #b.
      Pay(S, H1, H2, H3, R, c0, c1, c2, c3) @ #p
    & Redeemed(c0) @ #a
    & Refunded(c2) @ #b
    & not (Ex #k. Compromised(H2) @ #k)
    & not (Ex #k. Compromised(S) @ #k)"
\end{spthy}
Two control lemmas show that the PTLC result is not vacuous.
\code{Control_Refund_Of_H2_Outgoing} (exists-trace, 6 steps): H2's outgoing HTLC can be refunded
while H2 is honest, so the refund in the premise is not excluded by the model itself.
\code{Control_Wormhole_If_H2_Compromised} (exists-trace, 18 steps): once H2 itself is
compromised, the bypass is reachable with S still uncompromised, so the safety proof depends on
H2 being honest.
Sanity lemmas in both models show that a payment completes end to end without any compromise
(\code{Honest_Payment_Completes}: 10 steps in PTLC, 10 in HTLC) and with both neighbours H1 and
H3 of H2 compromised (\code{Colluders_Present_Payment_Completes}: 14 steps in PTLC, 8 in HTLC).
All eight lemmas are verified by \texttt{make aux} (PTLC 5/5, HTLC 3/3).

\paragraph{Limitations}
The two models are separate from \code{lightning.spthy}: a fixed route with three
intermediaries and one honest hop under attack, no channel layer, abstract timing, and an honest
\code{Sender_Setup} that chooses all locks and tweaks consistently; the safety lemma also
assumes that the sender is never compromised. The hash chain abstracts the group operation of
the adaptor-signature construction: it captures that a hop's incoming secret requires that hop's
tweak, but not the algebra of~\cite{malavolta2019} (an XOR encoding exhausted memory in
\tamarin). We therefore claim the contrast between the two lock designs on this route, not a
PTLC result for routes of any length.
